\documentclass[a4paper]{article}
\usepackage{amsmath}
\usepackage{amssymb}
\usepackage{geometry}
\usepackage{hyperref}
\usepackage{booktabs}
\usepackage{enumerate}
\geometry{margin=1in}

\begin{document}

\title{\textit{Introduction to Econometrics}\\Assignment 5}
\author{Richard Harnisch, s5238366}
\date{\today}
\maketitle

% answer environment
\makeatletter
\newcommand{\answer}[1]{%
    \par\medskip
    \noindent
    \hspace*{-\@totalleftmargin}%
    \fbox{%
        \parbox{\dimexpr\textwidth-2\fboxsep-2\fboxrule}{#1}%
    }%
    \par\medskip
}
\makeatother

\noindent\textit{TeX code available under\\ \url{https://github.com/richardharnisch/intro-to-econometrics}}

\begin{enumerate}

    \item Consider the model $y_i = \beta z_i + e_i$ where $z_i$ is a scalar and
          $E(z_i e_i) \neq 0$.  The dummy variable $D_i \in \{0,1\}$ satisfies
          $E(z_i D_i) \neq 0$.

          \begin{enumerate}

              \item What condition must $D_i$ satisfy with respect to $e_i$ to be a valid
                    instrument?

                    \answer{
                        $D_i$ must satisfy the \textbf{exogeneity} (exclusion restriction)
                        condition:
                        \[
                            E(D_i e_i) = 0.
                        \]
                        Together with the relevance condition $E(z_i D_i) \neq 0$ already given,
                        this makes $D_i$ a valid instrument for $z_i$.
                    }

              \item Using this condition, express $\beta$ as a function of
                    $E(D_i y_i)$ and $E(D_i z_i)$.

                    \answer{
                        Start from the model $y_i = \beta z_i + e_i$.  Multiply both sides by
                        $D_i$ and take expectations:
                        \[
                            E(D_i y_i) = \beta\, E(D_i z_i) + E(D_i e_i).
                        \]
                        Applying the exogeneity condition $E(D_i e_i) = 0$:
                        \[
                            E(D_i y_i) = \beta\, E(D_i z_i).
                        \]
                        Since $E(D_i z_i) \neq 0$ (relevance), we can solve for $\beta$:
                        \[
                            \beta = \frac{E(D_i y_i)}{E(D_i z_i)}.
                        \]
                    }

              \item Show that this expression can be rewritten solely as a function of
                    $E(y_i \mid D_i = 1)$ and $E(z_i \mid D_i = 1)$.

                    \answer{
                        Since $D_i \in \{0,1\}$, the Law of Total Expectation gives
                        \[
                            E(D_i y_i)
                            = E(D_i y_i \mid D_i=1)\,P(D_i=1)
                            + E(D_i y_i \mid D_i=0)\,P(D_i=0).
                        \]
                        When $D_i=1$ the product $D_i y_i = y_i$; when $D_i=0$ the product is
                        zero.  Therefore:
                        \[
                            E(D_i y_i) = E(y_i \mid D_i=1)\cdot P(D_i=1).
                        \]
                        Identically for the denominator:
                        \[
                            E(D_i z_i) = E(z_i \mid D_i=1)\cdot P(D_i=1).
                        \]
                        Substituting into the expression from (b), the factor $P(D_i=1) > 0$
                        cancels (it is positive because $E(D_i z_i) \neq 0$ requires
                        $P(D_i=1) \neq 0$):
                        \[
                            \beta
                            = \frac{E(y_i \mid D_i=1)\cdot P(D_i=1)}
                            {E(z_i \mid D_i=1)\cdot P(D_i=1)}
                            = \frac{E(y_i \mid D_i = 1)}{E(z_i \mid D_i = 1)}.
                        \]
                    }

          \end{enumerate}

    \item We want to estimate the causal effect of receiving disability insurance (DI)
          benefits on hours worked.

          \begin{enumerate}

              \item What concern would you have about estimating
                    $\text{hours worked}_i = \beta_0 + \beta_1 DI_i + \varepsilon_i$
                    by OLS?  How would the estimated $\hat{\beta}_1$ compare to the true
                    causal effect?

                    \answer{
                        \textbf{Endogeneity / selection bias.}  The indicator $DI_i$ is not
                        randomly assigned: individuals with more severe disabilities or worse
                        underlying health are both (i) more likely to receive DI benefits
                        ($DI_i = 1$) and (ii) less likely to work regardless of DI receipt.
                        Unobserved health severity is therefore captured in the error
                        $\varepsilon_i$, and it is negatively correlated with hours worked
                        while positively correlated with $DI_i$.  Formally,
                        $\mathrm{Cov}(DI_i, \varepsilon_i) < 0$, so OLS is inconsistent.

                        The OLS estimator conflates the causal effect of DI with the pre-existing
                        lower work capacity of DI recipients.  Consequently $\hat{\beta}_1$ will
                        be \textbf{more negative} (larger in absolute value) than the true causal
                        effect: OLS overestimates how much DI benefits reduce hours worked.
                    }

              \item Propose a variable and an estimator that could consistently estimate the
                    causal effect of DI benefits on hours worked.  Justify your answer.

                    \answer{
                        \textbf{Instrument.}  Use a measure of \emph{judge leniency}, e.g.\
                        the historical DI-application acceptance rate of the judge assigned to
                        individual $i$.

                        \textbf{Estimator.}  Instrumental Variables (IV), or equivalently
                        Two-Stage Least Squares (2SLS).

                        \textbf{Justification.}
                        \begin{itemize}
                            \item \textit{Relevance}: Judges differ greatly in acceptance rates,
                                  so a more lenient judge directly raises the probability that
                                  individual $i$ receives DI benefits.  The first-stage
                                  correlation $E(\text{leniency}_i \cdot DI_i) \neq 0$ is
                                  satisfied.

                            \item \textit{Exogeneity / exclusion restriction}: DI applications
                                  are randomly allocated to judges, so the assigned judge is
                                  independent of the applicant's health severity, earning
                                  capacity, and all other determinants of hours worked.  Judge
                                  leniency therefore affects hours worked \emph{only through}
                                  the DI-receipt decision, i.e.\
                                  $E(\text{leniency}_i \cdot \varepsilon_i) = 0$.
                        \end{itemize}
                        Because both conditions hold, IV/2SLS using judge leniency as an
                        instrument consistently estimates the causal effect $\beta_1$.
                    }

          \end{enumerate}

\end{enumerate}

\end{document}
